\documentclass[11pt,a4paper]{article}

\usepackage[utf8]{inputenc}
\usepackage[T1]{fontenc}
\usepackage{lmodern}
\usepackage[english]{babel}
\usepackage{geometry}
\usepackage{graphicx}
\usepackage{float}
\usepackage{caption}
\usepackage{amsmath}
\usepackage{booktabs}
\usepackage{hyperref}
\usepackage{subcaption}
\usepackage{xcolor}

\geometry{margin=1in}
\raggedbottom
\emergencystretch=3em   % let long \texttt{} identifiers break the line gracefully

% Figures live in the repo's runs/ directory; this file sits in report/5/,
% so resolve \includegraphics paths relative to the repository root.
\graphicspath{{../../}}

% Graceful fallback for figures produced by runs that may not be pulled yet.
\newcommand{\figorbox}[2]{%
  \IfFileExists{../../#1}{\includegraphics[width=#2\textwidth]{#1}}{%
    \fbox{\parbox[c][3.5cm][c]{#2\textwidth}{\centering\itshape
    \textcolor{gray}{[figure not yet generated --- run the eval and pull it]}}}}}

\title{Empirical Analysis of Transformers on Graph Connectivity \\
       \large Part V: Is It Matrix Powering or DFS?}
\author{Edoardo Paccagnella}
\date{\today}

\begin{document}

\maketitle

\begin{abstract}
\noindent The earlier reports read the model's sharp distance-$3^L$ wall as evidence that
the standard transformer computes connectivity by \emph{matrix powering} ($A^{3^L}$). This
report asks a different question, suggested by the advisor: what if the model is instead
learning a \emph{depth-first search} (DFS) --- a sequential traversal that explores a
component node by node and runs out of budget after a bounded number of steps, rather than a
parallel computation bounded by shortest-path \emph{distance}? The two mechanisms make
opposite predictions on graphs that are \emph{short but large} --- in particular on two dense
cliques joined by a single bridge edge --- and on whether \emph{density} helps or hurts. This
report sets out the hypothesis, the predictions that separate it from matrix powering, and the
experiments we will run; results follow in later revisions.
\end{abstract}

\section{Introduction}
\label{sec:intro}

\subsection{What the first four reports established}
\label{sec:recap}

We start this report from a set of conclusions we now treat as established (each was
multi-seed and is documented in Reports~I--IV). They are the ground this report builds on, and
they are also the facts the DFS hypothesis must remain consistent with.

\begin{itemize}\setlength\itemsep{0.25em}
  \item \textbf{The data lever does not help the standard transformer generalise.} Restricting
        the training data to within-capacity graphs (diameter $\le 3^L$), which is \emph{proven}
        to help the disentangled transformer, does \emph{not} improve out-of-distribution
        generalisation of a \emph{standard} transformer. Reproducing the paper's $n=20$ setup
        over many seeds, the outcome is \textbf{seed-dominated}: whether a run finds the
        generalising solution is a coin-flip ($\sim 3/8$ seeds in either condition), and the
        clean ``phase transition'' the paper shows is, in our hands, a single lucky seed
        (Report~III, Part~I). The filter only speeds up convergence (Report~II).
  \item \textbf{There is a sharp wall at shortest-path distance $\approx 3^L$ ($=9$ for our
        $L=2$ models), and it is architectural, not a data gap.} The wall moves with depth in a
        causal depth sweep ($L=1\to d{=}3$, $L=2\to d{=}9$, $L=3\to$ the whole range;
        Report~III). It appears even on families that are \emph{in} the training stream and seen
        $\sim 10^8$ times (2chains, 1cycle at $n=40$: exactly $1.00$ up to $d=9$, then exactly
        $0.00$; Report~IV, Thread~1). So steering the data cannot move it.
  \item \textbf{Density is hard at inference --- but so far confounded with being
        out-of-distribution.} A graph denser than the (sparse) training set fails even when its
        diameter is tiny (Report~IV, Thread~1, density axis). We attributed this to the graph
        being unlike anything seen in training. The single sharpest instance is the
        \textbf{degree heuristic} on dense 2cliques, which is repaired \emph{only by the right
        data} (mixed training), not by the read-out or the loss (Report~IV, Thread~3).
  \item \textbf{The spectral gap is not a separable axis on natural graphs}
        (collinear with the diameter, $-0.82$); a real bottleneck does add difficulty, but only
        a purpose-built, confound-free probe isolates it (Report~IV, Thread~2).
  \item \textbf{A similarity read-out doubles the reach to $2\cdot 3^L$} (meet-in-the-middle of
        two $3^L$-hop neighbourhoods) and removes the over-connection (cut) error, \emph{but
        cannot cross a bottleneck} --- that limit belongs to the trunk (Report~IV, Thread~3).
  \item \textbf{Counting components is trivial} for the same trunk (1-chain vs 2-chain: $100\%$
        in $2{,}000$ steps), so the bottleneck of the connectivity task is not sensing global
        structure but producing correct long-range \emph{per-pair} reachability (Report~III).
\end{itemize}

\subsection{The new question: matrix powering or DFS?}
\label{sec:question}

Everything above was \emph{interpreted} through one algorithm: matrix powering. An $L$-layer
model that repeatedly multiplies by the adjacency can confirm a connection within $3^L$ hops,
which fits the distance-$9$ wall exactly. But that wall is also consistent with a \emph{second}
algorithm we never separated from it: a \textbf{depth-first search}. A transformer cannot
recurse without bound in a fixed number of layers, so any DFS it learns must be a
\emph{bounded} traversal --- it explores a component step by step and \emph{stops after a
limited number of steps}, leaving the rest of the component unreached. The advisor's hypothesis
is that this, not matrix powering, may be what the model is doing: ``maybe the algorithm stops
after a certain iteration instead of running to the end, so it never finishes learning the
graph.''

The distinction matters because the two algorithms are bounded by \emph{different} quantities,
and on most graphs those quantities move together (a long graph has both a large distance and
many nodes to traverse), which is exactly why the first four reports could not tell them apart.
The point of this report is to build graphs where they \emph{disagree}.

\section{Two hypotheses, and how to tell them apart}
\label{sec:hypotheses}

\paragraph{Matrix powering (the reading of Reports~I--IV).} Connectivity via $A^{k}$ with
self-loops: $i$ and $j$ are connected iff $(A^{k})_{ij}>0$ for $k$ large enough. An $L$-layer
model realises this up to $k=3^L$. This computation is \emph{parallel} (every pair is resolved
at once), \emph{symmetric} (no preferred node or direction), and bounded by the
shortest-path \emph{distance} between the two nodes. Crucially, a \emph{denser} graph has
\emph{shorter} distances, so under matrix powering density should make connectivity
\emph{easier}, not harder.

\paragraph{DFS (the new hypothesis).} Connectivity by traversal: start at a node, walk into
neighbours, mark everything reachable as one component. A bounded DFS visits at most some fixed
number of nodes/edges before its budget runs out. This computation is \emph{sequential}, has a
\emph{starting point} (so it can resolve part of a component and not the rest), and is bounded
by the \emph{amount of structure to traverse} (number of nodes and edges), \emph{not} by
distance. Under DFS a \emph{denser} graph is \emph{harder}: each node has more neighbours, so a
fixed budget covers fewer nodes before it is exhausted.

\paragraph{Where they disagree.} Table~\ref{tab:predict} lists the predictions that separate
the two. The cleanest single discriminator is a graph that is \emph{short in distance but large
in nodes to traverse}: two dense cliques joined by a single edge. There the bridge sits at
distance $\le 3$ (trivial for matrix powering, well inside the $9$-hop capacity), but reaching
the far clique requires traversing one whole dense clique first (expensive for a bounded DFS).

\begin{table}[H]\centering\footnotesize\setlength{\tabcolsep}{4pt}
\begin{tabular}{lll}
\toprule
Observation & Matrix powering predicts & DFS predicts \\
\midrule
What bounds a correct answer & shortest-path \emph{distance} ($\le 3^L$) & \emph{nodes/edges} traversed \\
Dense clique + 1 bridge edge & easy (bridge is $\le 3$ hops) & hard (must cross a whole clique) \\
Effect of density at inference & easier (shorter distances) & harder (more to explore) \\
Effect of density in training & easier / faster to learn & harder / slower to learn \\
Error pattern in $\hat R$ & symmetric, no preferred node & ordered, a start node fully resolved \\
Asymmetry between two cliques & none & one clique can be right, the other partial \\
\bottomrule
\end{tabular}
\caption{Predictions that separate the two hypotheses. The model studied throughout is the
standard RoBERTa-faithful transformer with the \emph{linear} read-out of Reports~III--IV
($L=2$, single head, $d_{\mathrm{model}}=512$), unless stated otherwise.}
\label{tab:predict}
\end{table}

A note on consistency with prior results. The DFS reading does \emph{not} overturn the $3^L$
wall: on the sparse, low-degree families of Reports~III--IV (chains, cycles, path unions) every
node has $\approx 2$ neighbours, so ``distance'' and ``nodes traversed along the path'' are the
same number, and a step-bounded DFS and a distance-bounded matrix power are
\emph{indistinguishable} --- both wall at $9$. That is precisely why those experiments could not
separate the two, and why the discriminating experiments below all break that coincidence
(dense local structure, or density swept on its own).

\section{The experiments we will run}
\label{sec:plan}

This section records the plan so the logic is legible to a future reader; results will be added
as the runs complete. All experiments use the standard RoBERTa-faithful transformer with the
linear read-out (the ``base'' model of Reports~III--IV), trained at $n=20$ and $n=40$, multi-seed,
unless noted.

\subsection{Core test: barbell (one bridge) vs.\ two cliques}
\label{sec:plan-barbell}

\paragraph{The two labels.} On $n$ nodes, split into a left half and a right half of $n/2$.
\begin{itemize}\setlength\itemsep{0.1em}
  \item \textbf{Label A --- bridged cliques (one component).} Each half is a clique; a
        \emph{single} edge joins one node of the left clique to one node of the right. The
        connectivity target $R$ is the all-ones matrix.
  \item \textbf{Label B --- two cliques (two components).} The same two cliques, with
        \emph{no} bridge. $R$ is block-diagonal (left block and right block all-ones,
        cross-block all-zeros).
\end{itemize}
The two graphs differ by exactly one edge, and that edge creates a connection at distance
$\le 3$. (This is \emph{not} the existing \texttt{generate\_barbell\_graph}, which joins
$\sim n/3$ cliques through a long middle path; we added two new generators,
\texttt{generate\_bridged\_cliques\_graph} and \texttt{generate\_split\_cliques\_graph}, both
taking a \texttt{clique\_size} so we can sweep it.) We run this two ways:

\emph{(a) Eval the base connectivity model.} The RoBERTa, linear read-out model of
Reports~III--IV, both ER-trained and mixed-trained, at $n=20$ and $n=40$, evaluated on the
connectivity matrix --- what the base model already does. (2cliques is in the mixed training
stream; the bridged version is held out, so it is the clean test.)

\emph{(b) A dedicated, trainable classifier.} A graph-level binary classifier
(\texttt{GraphBinaryClassifier}: the same standard trunk, mean-pool $+$ one logit) trained from
scratch on a 50/50 bridged/split stream, to ask whether the architecture can be \emph{taught}
to tell the two labels apart at all, and how that depends on the clique size. We train one model
with the clique size drawn at random per sample and read accuracy \emph{by} clique size (flat
$=$ matrix powering; falling at large $c$ $=$ a bounded-traversal limit), plus a few fixed-$c$
runs to time convergence against density.

\paragraph{What we measure.}
\begin{enumerate}\setlength\itemsep{0.1em}
  \item \textbf{Can the base model tell A from B?} Both as the full connectivity matrix
        (does $\hat R$ contain the cross-block connection for A and not for B?) and, as a
        cleaner readout, as a graph-level discrimination of the two labels. Matrix powering
        $\Rightarrow$ easy at any clique size; DFS $\Rightarrow$ fails, and fails worse as the
        cliques grow denser/larger.
  \item \textbf{Where in $\hat R$ is the error (the ``stopped exploring'' signature)?} For
        label A, split the matrix into the left-clique block, the right-clique block, and the
        cross block, and report exact-match per block. The DFS signature is an \emph{asymmetry}:
        the clique containing the (implicit) start node is resolved (block $\approx 1.00$) while
        the cross block --- the far clique's connection to the near one --- is only partially
        right ($\approx 0.5$), because the traversal crossed the bridge but ran out of budget
        inside the far clique. Matrix powering predicts no such asymmetry.
  \item \textbf{Scaling with clique size.} Sweep $n/2 \in \{\dots\}$ (clique size). Under matrix
        powering the cross-block accuracy is flat in clique size (distance stays $\le 3$); under
        DFS it falls as the cliques grow (more to traverse before reaching / finishing the
        bridge).
\end{enumerate}

\subsection{Re-reading the existing barbell results for the same signature}
\label{sec:plan-reuse}

Report~IV already evaluated the (path-bridge) barbell per distance and found the mixed model
collapses at the bridge while getting within-clique pairs right. We will re-open those
predictions at the \emph{block} level (left clique / right clique / cross), to check for the
same asymmetry --- whether one side is resolved and the other is not --- without training
anything new. This is the cheap confirmation: the runs exist, we only add the per-block
breakdown.

\subsection{Density in optimisation, not just inference}
\label{sec:plan-density}

Reports~II--IV studied density (and diameter) at \emph{inference}. The DFS hypothesis makes a
prediction about \emph{training}: if the model learns a traversal whose cost grows with the
amount of structure, then \emph{denser} training graphs should make optimisation \emph{harder /
slower}, whereas matrix powering (shorter distances when dense) predicts the opposite. We will
sweep the training density (e.g.\ ER edge probability, or mean degree) at fixed $n$ and measure
convergence speed and final in-distribution accuracy, and the OOD transfer of each. We separate
this from the diameter sweep of Report~II (which varied distance, and found tighter $=$ faster):
here we hold the regime and vary \emph{density} on its own, the quantity DFS is sensitive to and
matrix powering is not.

\subsection{Reference algorithms: a DFS oracle and a matrix-power oracle}
\label{sec:plan-oracle}

To make ``looks like DFS'' a measurement rather than a story, we implement reference
connectivity-matrix predictors (\texttt{dfs\_oracle.py}) and compare the \emph{transformer's
error pattern} to each:
\begin{itemize}\setlength\itemsep{0.1em}
  \item a \textbf{matrix-power oracle}: $\mathbf{1}[(A+I)^{3^L}>0]$, the clean distance-bounded
        predictor (the cut oracle of Report~IV). On the bridged cliques it connects the two
        sides at \emph{any} clique size (the bridge is $\le 3$ hops), so its cross-block curve
        is flat at $1$.
  \item a \textbf{bounded-traversal oracle}: from each start node, explore at most a fixed
        \texttt{budget} of nodes and call them connected. The free parameter is the budget,
        i.e.\ the number of nodes the traversal covers before it stops.
\end{itemize}

\paragraph{A subtlety we found while building the oracle (recorded for honesty).} ``DFS'' is the
advisor's word for ``a traversal that stops early'', but the \emph{order} of exploration matters
on this graph. A literal depth-first search \emph{dives}: from a node in the near clique it
reaches the bridge endpoint within a couple of steps and crosses to the far clique almost
immediately, so a budget-bounded \emph{DFS} does \emph{not} get stuck in the near clique (we
checked: its cross-block rate is high). The clean ``gets stuck in the near clique'' behaviour is
the budget-bounded \emph{breadth}-first traversal, which first fills the whole near clique (all
$c-1$ neighbours at distance $1$) and only crosses the bridge once the budget exceeds the
near-clique size $c$ --- so it fails the cross pairs for budget $\le c$, sharply, however few
hops away the far clique is. We therefore use the \textbf{visit-bounded BFS} (a ball of the $k$
nearest nodes) as the primary ``bounded-traversal'' reference and keep the bounded DFS as a
second, partial-coverage reference. Both are visit-count-bounded, which is the property that
distinguishes them from the distance-bounded matrix power; the BFS is simply the cleaner
discriminator for two dense cliques on a single bridge.

We score, per family, which oracle the model's mistakes resemble --- agreement on the
\emph{cross-block} entries where the two oracles disagree, not just overall accuracy --- and how
that agreement moves with the clique size and with the budget. If the model's cross-block
accuracy falls with the clique size and tracks a small-budget traversal where it parts from the
matrix-power oracle, that is direct evidence for the bounded-traversal (DFS-like) reading; if it
stays flat at $1$ like the matrix-power oracle, that is evidence for matrix powering.

\paragraph{Implementation (this session).} New code: the two generators in \texttt{data.py};
\texttt{dfs\_oracle.py} (matrix-power, bounded-BFS, bounded-DFS); \texttt{eval\_bridged\_cliques.py}
(discrimination, per-block, asymmetry, clique-size sweep, oracle agreement; eval-only) with
\texttt{scripts/eval\_bridged\_cliques.sbatch} over the base checkpoints; the dedicated
classifier \texttt{experiments2/train\_bridged\_classifier.py} with
\texttt{scripts/bridged\_classifier.sbatch}; the density-in-training sweep reuses
\texttt{train\_families\_n20.py} (it already takes \texttt{--p}/\texttt{--n\_nodes}) via
\texttt{scripts/density\_sweep.sbatch}; plotting in \texttt{plot\_bridged\_cliques.py} and
\texttt{plot\_density\_sweep.py}. Outputs land under \texttt{runs/report5/}.

\section{Results}
\label{sec:results}

The runs of Section~\ref{sec:plan} are completing in order. We report two: the dedicated
trainable classifier (Section~\ref{sec:plan-barbell}, part~(b)), in
Section~\ref{sec:res-clf}; and --- the decisive one --- the base connectivity model evaluated on
the bridged cliques (Section~\ref{sec:plan-barbell}, part~(a), with the per-block,
clique-size-sweep and oracle-agreement reads), in Section~\ref{sec:res-base}. The
training-density sweep and the final matrix-powering-vs-DFS verdict will be added as they land.

\subsection{The dedicated classifier: the bridge is detectable at every clique size}
\label{sec:res-clf}

\paragraph{Setup recap.} We train the binary classifier of Section~\ref{sec:plan-barbell}(b)
(the standard minimal trunk --- $L=2$, single head, $d_{\mathrm{model}}=512$, normalised-ReLU ---
with mean-pool $+$ one logit) from scratch at $n=20$, $200{,}000$ steps. The data is an
\emph{online stream}: there is no fixed training set, each graph is freshly generated at every
step, with the bridged/split label drawn $50/50$. A graph of ``clique size $c$'' is two cliques of
$c$ nodes (so $2c$ of the $n=20$ nodes are used, the remaining $20-2c$ left isolated), bridged or
not; $c$ can range over $\{2,3,\dots,10\}$ ($c=10$ fills the whole canvas, no padding).

The two \textbf{regimes} differ only in how $c$ is chosen for each graph of the stream, and they
ask two different questions:
\begin{itemize}\setlength\itemsep{0.1em}
  \item \textbf{Random $c$ (the discriminator).} For \emph{every} training graph, $c$ is drawn
        uniformly from $\{2,\dots,10\}$, so a \emph{single} model sees all clique sizes mixed
        together (each $\approx 1/9$ of the stream). At evaluation we read accuracy \emph{separately
        per clique size}: does the \emph{same} model get worse as the cliques grow? Flat in $c$ is
        the matrix-powering reading, falling at large $c$ is the bounded-traversal (DFS-like)
        reading. Four seeds.
  \item \textbf{Fixed $c$ (convergence vs.\ density).} The whole stream uses a single clique size;
        we train a separate model on $c=3$, on $c=6$ and on $c=10$ (two seeds each). Here the
        question is about \emph{optimisation}, not capacity: is a denser clique slower or harder to
        learn? We compare convergence speed and loss across the three densities.
\end{itemize}
The balanced test set has $200$ graphs per label per clique size (covering $c=2,\dots,10$ in the
random regime, the single $c$ in the fixed regime).

\paragraph{Final accuracy is flat at $1.0$ across every clique size.} By the end of training all
four random-$c$ seeds reach $\approx1.00$ accuracy at \emph{every} clique size $c=2,\dots,10$
(the largest pair of cliques fills the whole $n=20$ canvas); the worst single cell over all
seeds and sizes is $0.995$. The fixed-$c$ runs tell the same story: $c=3$, $6$ and $10$ all reach
$1.00$, and in fact already at the first evaluation ($2{,}000$ steps). So there is \emph{no
permanent accuracy wall} in the clique size: the architecture can be taught to separate the two
labels however large the cliques get, which is the matrix-powering-consistent outcome and rules
out the strongest form of the DFS hypothesis (that a bridge behind a large clique is simply
\emph{undetectable}).

\paragraph{But the DFS difficulty ordering does show up --- in the learning dynamics.} The flat
final curve hides a transient that is exactly the shape DFS predicts. Figure~\ref{fig:clf} plots,
for each seed, the accuracy of every clique size as a function of the training step (one line per
size, colour $=$ size): the small cliques shoot up to $1.0$ almost immediately, the larger ones
only later, so the curves fan out by size before all converging --- the larger, denser cliques are
the \emph{last} to be mastered. Read as a snapshot early in training
(Table~\ref{tab:clf_early}, step $2{,}000$): seed $1000$ is $1.00$ up to $c=6$ but only $0.65$ at
$c=8$ and $0.55$ ($\approx$ chance) at $c=10$; seed $3000$ is milder ($0.87$ at $c=10$), seed
$4000$ milder still, seed $2000$ already flat. This is a bounded-traversal-like ordering that the
model overcomes with more steps, not a hard limit. The fixed-$c$ runs corroborate the optimisation
side (Figure~\ref{fig:clf_loss}): although accuracy reaches $1.0$ at every $c$, the early
($2{,}000$-step) training \emph{loss} rises with the clique size ($\approx0.05$ at $c=3$,
$\approx0.17$ at $c=6$, $\approx0.35$ at $c=10$), so a denser clique is genuinely a bit harder to
fit --- the ``density in optimisation'' effect DFS predicts --- just not hard enough to delay
reaching $100\%$.

\begin{table}[H]\centering\small
\begin{tabular}{lccccccccc}
\toprule
seed & $c2$ & $c3$ & $c4$ & $c5$ & $c6$ & $c7$ & $c8$ & $c9$ & $c10$ \\
\midrule
1000 & 1.00 & 1.00 & 1.00 & 0.97 & 0.99 & 0.91 & \textbf{0.65} & \textbf{0.59} & \textbf{0.55} \\
2000 & 1.00 & 1.00 & 1.00 & 1.00 & 1.00 & 1.00 & 1.00 & 1.00 & 1.00 \\
3000 & 1.00 & 1.00 & 0.99 & 0.99 & 0.99 & 0.93 & 0.95 & 0.92 & \textbf{0.87} \\
4000 & 1.00 & 1.00 & 1.00 & 1.00 & 0.99 & 1.00 & 0.98 & 0.99 & 0.97 \\
\bottomrule
\end{tabular}
\caption{\emph{Model:} the bridged-vs-split binary classifier --- minimal/A.1-style trunk ($L=2$,
single head, $d_{\mathrm{model}}=512$, normalised-ReLU) with mean-pool $+$ one logit --- trained
from scratch at $n=20$ ($200{,}000$ steps) on the online $50/50$ bridged/split stream with the
clique size drawn uniformly per sample (\emph{random-$c$} regime), four seeds.
\emph{Test set:} balanced, $200$ graphs per label per clique size $c=2,\dots,10$.
Classification accuracy \textbf{by clique size, early in training (step $2{,}000$)}; bold marks the
cells near chance. The larger cliques are learned last (seed-dependently); by the end of training
every cell is $\approx1.00$ (the trajectories all converge, Figure~\ref{fig:clf}), so this is a
transient ordering, not a wall.}
\label{tab:clf_early}
\end{table}

\begin{figure}[H]
    \centering
    \figorbox{runs/report5/report5_figs/classifier_trajectory_n20.png}{0.92}
    \caption{\emph{Model:} the bridged-vs-split classifier (minimal trunk $L=2$, single head,
    $d_{\mathrm{model}}=512$, mean-pool $+$ one logit), trained from scratch at $n=20$ on the online
    $50/50$ bridged/split stream, \emph{random-$c$} regime (clique size drawn uniformly
    $\{2,\dots,10\}$ per sample), one panel per seed. \emph{Test set:} balanced, $200$ graphs per
    label per clique size. \textbf{Accuracy vs training step (log axis), one line per clique size}
    (colour $=$ clique size $c$). Small cliques (dark) reach $1.0$ almost at once, the larger ones
    (yellow) later, and every size converges to $1.0$ --- the difficulty ordering is in \emph{when}
    a size is learned, not in whether it is. (At $n=20$, $c\le10$, the spread is small and partly
    hidden by mid-run noise; it is much clearer at $n=40$, Figure~\ref{fig:clf_n40}.)}
    \label{fig:clf}
\end{figure}

\paragraph{Pushing the range: $n=40$, cliques up to $c=20$.} At $n=20$ the largest clique is
$c=10$, so the ``flat at $1.0$'' conclusion is only tested that far; a bounded-traversal limit, if
real, might only bite at larger cliques. We therefore re-ran the whole classifier at $n=40$, where
the clique size reaches $c=20$ (two $20$-cliques fill the entire canvas), with fixed sizes
$c\in\{5,10,15,20\}$. The conclusion is the \emph{same, only stronger}. The final accuracy is
\textbf{exactly $1.0$ at every clique size $c=2,\dots,20$, for all four seeds} (no dips at all,
cleaner than $n=20$): doubling the range does not produce a wall, so the bridge stays detectable
behind even a $20$-node dense clique --- the matrix-powering-consistent reading, and a sharper
refutation of the strong DFS hypothesis.

\paragraph{The bounded-traversal cost reappears --- as an optimisation cost, more clearly than at
$n=20$.} With the larger range the difficulty ordering is now unmistakable. The per-size
trajectories (Figure~\ref{fig:clf_n40}) fan out widely: small cliques are learned almost at once,
the largest only after tens of thousands of steps. As a snapshot at step $2{,}000$
(Table~\ref{tab:clf_early_n40}) accuracy falls \emph{monotonically} with the clique size: seed
$4000$ is $1.00$ at $c=2$, $0.77$ at $c=8$ and flat at chance ($\approx0.50$) from $c\ge11$, with
seed $3000$ similar and seeds $1000$/$2000$ slower still --- the larger the clique, the later it is
mastered. The fixed-$c$ runs make the optimisation side explicit and sharp
(Figure~\ref{fig:clf_loss}): the step-$2{,}000$ training loss climbs from $\approx0.1$ at $c=5$ to
$\approx0.69$ ($=\ln 2$, i.e.\ \emph{still at chance}) at $c=15$ and $c=20$, and the step at which
accuracy first reaches $0.99$ grows from $2{,}000$ (at $c=5,10$) to $4{,}000$--$6{,}000$ (at
$c=15,20$). So a
denser/larger clique is genuinely harder and slower to learn --- exactly the ``density in
optimisation'' prediction of the DFS reading --- but, given enough steps, never an accuracy
ceiling. The bounded-traversal difficulty shows up as a \emph{training cost}, not a capability
limit. (At $n=20$, where $c\le10$, this was barely visible: every fixed-$c$ run was already at
$1.0$ by the first evaluation.)

\begin{table}[H]\centering\small
\begin{tabular}{lccccccc}
\toprule
seed & $c2$ & $c5$ & $c8$ & $c11$ & $c14$ & $c17$ & $c20$ \\
\midrule
1000 & 0.62 & 0.69 & 0.67 & 0.66 & 0.59 & 0.56 & \textbf{0.49} \\
2000 & \textbf{0.50} & 0.55 & \textbf{0.47} & \textbf{0.50} & 0.52 & \textbf{0.50} & \textbf{0.50} \\
3000 & 0.98 & 0.76 & \textbf{0.51} & \textbf{0.51} & \textbf{0.50} & \textbf{0.50} & \textbf{0.50} \\
4000 & 1.00 & 0.88 & 0.77 & \textbf{0.52} & \textbf{0.52} & \textbf{0.50} & \textbf{0.50} \\
\bottomrule
\end{tabular}
\caption{\emph{Model:} the bridged-vs-split classifier (minimal trunk $L=2$, single head,
$d_{\mathrm{model}}=512$, mean-pool $+$ one logit), trained from scratch at \textbf{$n=40$} on the
online $50/50$ bridged/split stream, \emph{random-$c$} regime (clique size drawn uniformly
$\{2,\dots,20\}$ per sample), four seeds. \emph{Test set:} balanced, $200$ graphs per label per
clique size. Classification accuracy \textbf{by clique size, early in training (step $2{,}000$)}
(cols thinned to $c=2,5,\dots,20$); bold marks cells at chance ($\approx0.50$). The drop with
clique size is monotone and much deeper than at $n=20$ (Table~\ref{tab:clf_early}); by the end of
training every cell is exactly $1.00$ for every seed.}
\label{tab:clf_early_n40}
\end{table}

\begin{figure}[H]
    \centering
    \figorbox{runs/report5/report5_figs/classifier_trajectory_n40.png}{0.92}
    \caption{\emph{Model:} the bridged-vs-split classifier (minimal trunk $L=2$, single head,
    $d_{\mathrm{model}}=512$, mean-pool $+$ one logit), trained from scratch at \textbf{$n=40$} on
    the online $50/50$ bridged/split stream, \emph{random-$c$} regime (clique size drawn uniformly
    $\{2,\dots,20\}$ per sample), one panel per seed. \emph{Test set:} balanced, $200$ graphs per
    label per clique size. Same layout as Figure~\ref{fig:clf}: accuracy vs training step (log
    axis), one line per clique size (colour $=$ size). The fan is wide: small cliques (dark) reach
    $1.0$ almost at once, the largest (yellow) only after tens of thousands of steps, but
    \emph{every} size converges to $1.0$. The ordering is monotone in $c$ and far clearer than at
    $n=20$.}
    \label{fig:clf_n40}
\end{figure}

\begin{figure}[H]
    \centering
    \figorbox{runs/report5/report5_figs/classifier_fixedc_loss.png}{0.6}
    \caption{\emph{Model:} the bridged-vs-split classifier (minimal trunk $L=2$, single head,
    $d_{\mathrm{model}}=512$, mean-pool $+$ one logit), but a \emph{separate} model trained from
    scratch on each single \emph{fixed} clique size (\emph{fixed-$c$} regime: the online $50/50$
    bridged/split stream uses one $c$), $n=20$ ($c\in\{3,6,10\}$) and $n=40$ ($c\in\{5,10,15,20\}$),
    two seeds each. \emph{Plotted:} the \textbf{training} loss (on the stream, not a test metric) at
    step $2{,}000$ vs the clique size (markers $=$ the two seeds, line $=$ mean). The loss rises
    with the clique size and saturates at the chance value $\ln 2$ (dotted) for the largest $n=40$
    cliques --- a denser clique is genuinely slower to fit (``density in optimisation''). The
    held-out accuracy nonetheless reaches $1.0$ at every size with more steps, so this is a
    convergence-speed effect, not a capability wall.}
    \label{fig:clf_loss}
\end{figure}

\paragraph{Two caveats, stated plainly.} First, training is \emph{noisy} at both sizes: some seeds
suffer transient mid-run collapses to chance on the hardest cliques before settling (the step at
which all clique sizes stabilise above $0.99$ ranges from $\sim4{,}000$ to $\sim116{,}000$ across
the $n=40$ seeds), so the clique-size ordering is a tendency of the early dynamics, not a clean
monotone law --- the familiar seed lottery of Reports~III--IV, here at the level of which clique
sizes are mastered when. Second, and more important for the headline
question: this classifier is a \emph{weak} discriminator of matrix powering vs.\ DFS. The two
labels differ by exactly one edge, and that edge leaves a purely \emph{local} footprint --- the
two bridge endpoints have degree $c$ instead of $c-1$, and the total edge count differs by one ---
so the labels are separable without ever propagating connectivity across the cliques. A flat
$1.0$ therefore shows the bridge is \emph{detectable} at any clique size; it does \emph{not} show
that the model computes the full connectivity matrix by matrix powering. That stronger question
--- whether the base connectivity model actually fills $\hat R$ correctly across the bridge, and
whether its errors carry the bounded-traversal signature --- is what the base-model eval
(Section~\ref{sec:plan-barbell}(a)) and the oracle comparison (Section~\ref{sec:plan-oracle})
decide. We turn to it now, and it gives a sharper answer than the classifier.

\subsection{The base model: the single bridge is not propagated across a large clique}
\label{sec:res-base}

\paragraph{Setup recap.} This is the decisive test (Section~\ref{sec:plan-barbell}(a)). We take
the base connectivity model of Reports~III--IV --- the RoBERTa-faithful transformer with the
\emph{linear} read-out ($L=2$, single head, $d_{\mathrm{model}}=512$), the model that emits the
full $n\times n$ connectivity matrix $\hat R$ --- and feed it the bridged and split cliques
\emph{without any further training} (eval-only). We use both the ER-trained and the mixed-trained
checkpoints, at $n=20$ and $n=40$, four seeds each ($16$ checkpoints). The \textbf{mixed} model is
the clean test: dense cliques (2cliques, clique-blocks) are in its training stream, so it is
\emph{not} out-of-distribution on the cliques themselves; the \textbf{ER} model has only seen
sparse graphs, so on dense cliques it is in the degree-heuristic OOD regime of Report~IV
(within-clique accuracy already $<1$) and we read it only as a confounded cross-check, not as
evidence on its own (Report~IV, Thread~1). For each bridged graph we split $\hat R$ into the
left-clique block, the right-clique block and the cross block, and we sweep the clique size
$c$ (so $2c$ nodes are used, the rest left isolated). The bridge always sits at distance $\le 3$,
\emph{at every $c$} --- so matrix powering must score the cross block $1.0$ for all $c$, while a
visit-bounded traversal must lose it once one clique alone exhausts the budget.

\paragraph{What ``cross-block pair accuracy'' means (and why it is the quantity that matters).}
The model outputs an $n\times n$ matrix $\hat R$; entry $\hat R_{ij}$ is its guess for ``are nodes
$i$ and $j$ in the same component?'' A \emph{pair} is one such $(i,j)$ entry, and \emph{pair
accuracy} on a set of pairs is just the fraction of those entries the model gets right (predicted
$0/1$ equals the true $0/1$). We group the pairs into three blocks: \emph{within-A} (both nodes in
the left clique), \emph{within-B} (both in the right clique), and \emph{cross} (one node in each
clique). On a \emph{bridged} graph the two cliques are one component, so the true label of
\emph{every} cross pair is $1$ (connected); cross-block accuracy is therefore the fraction of A--B
pairs the model correctly calls connected, and $\approx 0.00$ means it calls essentially all of
them \emph{disconnected}. The key thing to keep separate: the bridge is a single \emph{edge}, and
it lives in the \emph{input} adjacency $A$ (the model can see it); the cross \emph{block} lives in
the \emph{output} $\hat R$ and has $c^2$ entries. Filling it correctly does not mean copying one
edge --- it means \emph{propagating} that one edge into ``and therefore all $c^2$ cross pairs are
connected.'' Seeing the edge and propagating it are different operations, and that distinction is
the whole story below.

\paragraph{Within each clique: perfect. Across the bridge: collapses with clique size.} The mixed
model resolves each clique \emph{perfectly} and reports two components on the split graph
\emph{perfectly} --- within-A, within-B and the whole split matrix are all exactly $1.00$, every
seed, both $n$ (Table~\ref{tab:base_blocks}). So the failure is \emph{not} that dense cliques are
hard (they are in-distribution; the model gets them right). The only thing the mixed model misses
is the \emph{cross} block of the bridged graph: at full clique size ($c=10$ at $n=20$, $c=20$ at
$n=40$) the cross-block accuracy is $\approx 0.00$ --- it sees the two bridged cliques as
\emph{two} components, exactly as if the bridge were not there. And the collapse is graded in the
clique size (Figure~\ref{fig:base_sweep}, Table~\ref{tab:base_sweep}): cross accuracy is flat at
$1.0$ for small cliques and falls to $0$ as the cliques grow ($n=20$: $1.0$ up to $c=6$, then
$0.27$ at $c=7$, $0.01$ at $c=10$; $n=40$: $1.0$ up to $c=6$, $0.74$ at $c=7$, $0.27$ at $c=8$,
$0$ from $c\simeq14$). The matrix-power oracle, on the same graphs, is flat at $1.0$ for every $c$
(the bridge is $\le 3$ hops). \textbf{The model departs from matrix powering exactly where a
bounded traversal would: once the near clique alone is too big to ``get through'' before reaching
the far one.}

\paragraph{The cutoff is an absolute node budget, not a fraction of the graph.} The collapse
happens at roughly the \emph{same absolute clique size} at $n=20$ (knee at $c\approx6$--$7$) and
$n=40$ (knee at $c\approx7$--$8$), even though $c=7$ is $70\%$ of the $n=20$ canvas but only $35\%$
of the $n=40$ one. A traversal that can cover about $6$--$7$ nodes before running out of budget
explains both. That is precisely the quantity the DFS hypothesis says bounds the model --- the
\emph{number of nodes traversed} --- and \emph{not} the shortest-path distance, which is $\le 3$
at every $c$ and would predict no collapse at all.

\paragraph{No asymmetry: it is a parallel, symmetric bounded \emph{ball}, not a single-start DFS.}
We had predicted (Table~\ref{tab:predict}) a \emph{directional} DFS signature --- one clique fully
resolved, the far one only partial, $\hat R_{ij}\neq\hat R_{ji}$. We do \emph{not} see it: both
cliques are resolved equally ($1.00$ vs $1.00$), the two cross directions fail equally
($\hat R_{AB}\approx\hat R_{BA}\approx 0$), and the directional-disagreement rate is $\approx 0$.
Stated honestly, the literal single-start-DFS picture is wrong here, and for a clear reason: a
transformer emits all of $\hat R$ \emph{in parallel}, so there is no single start node. The right
reference is the one we anticipated in Section~\ref{sec:plan-oracle} --- the
\emph{visit-bounded BFS ball}: every node resolves the cluster of nodes within its budget (its own
clique, which it gets exactly right), and the bridge into the far clique is dropped once the near
clique alone fills the budget.

\paragraph{How the oracle comparison works (and why the oracles are not ``always right'').} The two
oracles are \emph{not} scored for correctness --- they are two \emph{reference algorithms}, and we
ask which one the model's answers (mistakes included) look like. One of them, the \emph{matrix-power
oracle} $\mathbf 1[(A+I)^{9}>0]$, does happen to be the true connectivity, so on these graphs it is
right everywhere (cross block $=1$). The other, the \emph{visit-bounded BFS} with a fixed budget, is
\emph{deliberately} an early-stopping algorithm: it explores a ball of the nearest budget-many nodes
and calls only those connected, so on the bridged cliques it gets the cross block \emph{wrong} on
purpose (it never leaves the near clique). They are useful precisely because they \emph{disagree},
and exactly on the cross block. So the comparison is: take the model's $0/1$ for each pair and ask
what fraction matches each oracle's $0/1$. ``Overall agreement with the matrix-power oracle is only
$0.48$'' means the model's matrix matches the \emph{true} connectivity on only $48\%$ of pairs ---
which is just the within-block pairs; it disagrees with the truth on essentially the whole cross
block. The visit-bounded BFS, on the other hand, the model matches almost perfectly: at $n=20$,
with the budget set to one clique ($b=c=10$), the model agrees with the BFS oracle on $99\%$ of all
pairs, and on the contested cross entries --- the pairs where the two oracles split --- it sides
with the BFS oracle $\approx 99.6\%$ of the time and with matrix powering $0.4\%$
(Figure~\ref{fig:base_oracle}; identical across the four seeds, and the same direction at $n=40$,
where it reaches $\sim\!79\%$ BFS on the contested entries at $b=10$ and keeps climbing toward the
clique size). In words: the model is the early-stopping traversal almost exactly, and matrix
powering only on the parts where the two cannot be told apart.

\paragraph{The ER model is the confounded cross-check (do not read it alone).} The ER curves
(Figure~\ref{fig:base_sweep}, bottom) fall with $c$ too, but for a mixture of reasons we cannot
disentangle (Table~\ref{tab:base_er}): the ER model is already wrong \emph{within} the cliques
(degree-heuristic OOD, within-clique accuracy $0.0$--$1.0$ depending on the seed), and it is far
noisier seed-to-seed
(one $n{=}40$ seed collapses to $0$ everywhere; another rises back toward $1$ at large $c$ only
because the degree heuristic calls everything dense ``connected'' --- it then \emph{also} wrongly
calls the \emph{split} cliques connected, so its discrimination stays at chance). A falling ER
curve is therefore not clean evidence of bounded traversal --- density-OOD and bridge-distance
grow together (Report~IV, error noted in our handoff). This is exactly why the mixed model is the
test that counts, and there the within-clique blocks are perfect, so the cross collapse is
attributable to the bridge alone.

\paragraph{Why this does not contradict the classifier (cross block $\approx 0$ here, accuracy
$1.0$ there).} This is the natural objection, so it is worth spelling out: in
Section~\ref{sec:res-clf} the classifier reached $1.0$ at every clique size, and here the same kind
of model scores $0.00$ on the same bridged graph. There is no contradiction, because the two models
are asked \emph{different questions}. The classifier emits \emph{one bit} --- ``is there a bridge,
yes or no?'' --- and that bit is readable from a purely \emph{local} fingerprint of the bridge edge:
the two endpoints have degree $c$ instead of $c-1$, and the graph has one more edge in total. You
can tell bridged from split by counting, without ever working out who is connected to whom. (This is
the weak-discriminator caveat of Section~\ref{sec:res-clf}: detecting the edge is easy.) The base
model is asked the \emph{harder} question --- fill in all $c^2$ cross pairs of $\hat R$ --- and that
requires \emph{propagating} the single bridge edge into ``therefore every node on the left is
connected to every node on the right.'' It is this propagation, not the detection, that fails:
the model plainly \emph{sees} the bridge edge (it is right there in the input $A$, and the classifier
proves it is usable), but it does \emph{not} carry that connection across a clique of more than
$\sim 6$ nodes when it builds the connectivity matrix. So the user's intuition --- ``can't it just
read off the $1$ at the bridge position?'' --- is exactly what the classifier does; the base model's
job is the opposite of reading off one entry, it is spreading that one entry over the whole cross
block, and that is where the bounded traversal runs out. On the decisive read --- a
model that already handles the dense cliques, failing the bridge in a graded, absolute-node-budget
way that tracks a visit-bounded traversal and parts from the matrix-power oracle --- the base-model
eval comes down on the side of a \emph{bounded-traversal} mechanism, not distance-bounded matrix
powering.

\begin{table}[H]\centering\small
\begin{tabular}{llcccc}
\toprule
 & & \multicolumn{3}{c}{bridged graph (one component)} & split graph \\
\cmidrule(lr){3-5}\cmidrule(lr){6-6}
$n$ & seed & within-A & within-B & cross & exact-match \\
\midrule
\multicolumn{6}{l}{\emph{$n=20$, clique size $c=10$ (fills the canvas)}} \\
20 & 1000 & 1.00 & 1.00 & \textbf{0.00} & 1.00 \\
20 & 2000 & 1.00 & 1.00 & \textbf{0.01} & 1.00 \\
20 & 3000 & 1.00 & 1.00 & \textbf{0.01} & 1.00 \\
20 & 4000 & 1.00 & 1.00 & \textbf{0.01} & 1.00 \\
\multicolumn{6}{l}{\emph{$n=40$, clique size $c=20$ (fills the canvas)}} \\
40 & 1000 & 1.00 & 1.00 & \textbf{0.00} & 1.00 \\
40 & 2000 & 1.00 & 1.00 & \textbf{0.00} & 1.00 \\
40 & 3000 & 1.00 & 1.00 & \textbf{0.00} & 1.00 \\
40 & 4000 & 1.00 & 1.00 & \textbf{0.00} & 1.00 \\
\bottomrule
\end{tabular}
\caption{\emph{Model:} the base connectivity transformer --- RoBERTa-faithful trunk, \emph{linear}
read-out ($L=2$, single head, $d_{\mathrm{model}}=512$), emitting the full $n\times n$ matrix
$\hat R$ --- \textbf{mixed}-trained (the $9$-family online stream of Reports~III--IV, which includes
dense cliques), four seeds, $n=20$ and $n=40$; \emph{eval-only}, no further training.
\emph{Test set:} bridged cliques and split cliques at the full clique size ($2{,}000$ graphs each).
\emph{Metric:} per-block pair accuracy (within each clique, and the cross block A$\leftrightarrow$B)
on the bridged graph, and exact-match on the split graph; per seed. The within-clique blocks and
the entire split matrix are perfect; only the \emph{cross} block (the single bridge) is missed
(bold), at every seed. The two cross directions agree (no asymmetry), and the discrimination
accuracy is $0.50$ (the model calls bridged and split alike --- both ``two components''). The ER
model is omitted here (within-clique accuracy $<1$, confounded; see text).}
\label{tab:base_blocks}
\end{table}

\begin{figure}[H]
    \centering
    \figorbox{runs/report5/report5_figs/base_bridged_cross_sweep.png}{0.95}
    \caption{\emph{Model:} the base connectivity transformer (RoBERTa trunk, linear read-out,
    $L=2$, single head, $d_{\mathrm{model}}=512$, full-matrix $\hat R$), eval-only; \emph{top row}
    \textbf{mixed}-trained (clean: dense cliques are in training), \emph{bottom row} \textbf{ER}-trained
    (confounded: degree-heuristic OOD on dense cliques); $n=20$ (left), $n=40$ (right); four seeds
    each (thin lines $=$ seeds, thick $=$ mean). \emph{Test set:} bridged cliques, $\ge 400$ graphs
    per clique size. \emph{Metric:} cross-block pair accuracy (does $\hat R$ connect the two
    cliques?) vs clique size $c$. The bridge sits at distance $\le 3$ at \emph{every} $c$, so the
    matrix-power oracle (green dashed) is flat at $1$; the mixed model instead collapses from $1$ to
    $0$ once the near clique exceeds $\sim 6$--$7$ nodes --- at roughly the same \emph{absolute}
    clique size for both $n$ --- the visit-bounded-traversal prediction. The ER curves fall too but
    are noisy and confounded (see text), and are shown only as a cross-check.}
    \label{fig:base_sweep}
\end{figure}

\begin{table}[H]\centering\small\setlength{\tabcolsep}{4.5pt}
\begin{tabular}{lccccccccc}
\toprule
clique size $c$ & 2 & 3 & 4 & 5 & 6 & 7 & 8 & 10 & 14 \\
\midrule
model, $n=20$ (mean) & 1.00 & 1.00 & 1.00 & 1.00 & 0.95 & 0.27 & 0.13 & 0.01 & --- \\
model, $n=40$ (mean) & 0.89$^\dagger$ & 1.00 & 1.00 & 1.00 & 0.95 & 0.74 & 0.27 & 0.09 & 0.00 \\
matrix-power oracle  & 1.00 & 1.00 & 1.00 & 1.00 & 1.00 & 1.00 & 1.00 & 1.00 & 1.00 \\
\bottomrule
\end{tabular}
\caption{\emph{Model:} the \textbf{mixed}-trained base connectivity transformer (RoBERTa trunk,
linear read-out, $L=2$, single head, $d_{\mathrm{model}}=512$, full-matrix $\hat R$), eval-only,
four seeds, $n=20$ and $n=40$. \emph{Test set:} bridged cliques, $\ge 400$ graphs per clique size.
\emph{Metric:} cross-block pair accuracy on the bridged graph, \emph{mean over the four seeds}, by
clique size $c$ (selected columns; $c$ runs to $10$ at $n=20$ and to $20$ at $n=40$, all $0.00$
from $c=14$ on). The matrix-power oracle stays flat at $1$ because the bridge is always $\le 3$
hops; the model collapses with $c$. $^\dagger$the $n=40$, $c=2$ mean is pulled down by one noisy
seed ($0.58$; the other three are $\approx 1.00$).}
\label{tab:base_sweep}
\end{table}

\begin{figure}[H]
    \centering
    \figorbox{runs/report5/report5_figs/base_bridged_oracle_follow.png}{0.95}
    \caption{\emph{Model:} the \textbf{mixed}-trained base connectivity transformer (RoBERTa trunk,
    linear read-out, $L=2$, single head, $d_{\mathrm{model}}=512$, full-matrix $\hat R$), eval-only,
    four seeds; $n=20$ (left), $n=40$ (right). \emph{Test set:} bridged cliques at full clique size,
    $500$ graphs. \emph{Metric:} restricted to the cross pairs where the matrix-power oracle and the
    visit-bounded-BFS oracle \emph{disagree} (the bridge pairs), the fraction the model matches to
    each oracle, as the BFS budget $b$ varies. The model tracks the early-stopping BFS, not matrix
    powering: as $b$ approaches one clique (dashed line) the model matches BFS on $\to 1.0$ of the
    contested pairs and matrix powering on $\to 0$. This is the same fact as the cross collapse,
    read against the two reference algorithms.}
    \label{fig:base_oracle}
\end{figure}

\begin{table}[H]\centering\small
\begin{tabular}{llcccc}
\toprule
$n$ & seed & within-A & within-B & cross & discrimination \\
\midrule
20 & 1000 & 1.00 & 1.00 & 0.09 & 0.50 \\
20 & 2000 & 0.26 & 0.25 & 0.18 & 0.50 \\
20 & 3000 & 0.62 & 0.61 & 0.04 & 0.50 \\
20 & 4000 & 0.69 & 0.70 & 0.66 & 0.65 \\
40 & 1000 & 0.30 & 0.30 & 0.24 & 0.50 \\
40 & 2000 & 0.00 & 0.00 & 0.00 & 0.50 \\
40 & 3000 & 0.96 & 0.96 & 0.93 & 0.50 \\
40 & 4000 & 0.02 & 0.02 & 0.00 & 0.50 \\
\bottomrule
\end{tabular}
\caption{\emph{Model:} the \textbf{ER}-trained base connectivity transformer (RoBERTa trunk, linear
read-out, $L=2$, single head, $d_{\mathrm{model}}=512$, full-matrix $\hat R$) --- the
\emph{confounded} cross-check --- eval-only, four seeds, $n=20$ and $n=40$. \emph{Test set:} bridged
cliques at full clique size ($c=10$ at $n=20$, $c=20$ at $n=40$), $2{,}000$ graphs. \emph{Metric:}
per-block pair accuracy on the bridged graph (within each clique, and the cross block), and the
graph-level discrimination accuracy. Contrast with the clean mixed model (Table~\ref{tab:base_blocks},
where within-A and within-B are a flat $1.00$): here the within-clique accuracy is all over the place
($0.00$ to $1.00$) and swings wildly by seed, because the ER model has never seen a dense clique and
is in the degree-heuristic OOD regime of Report~IV --- it is already wrong \emph{inside} the cliques,
so a low cross accuracy cannot be blamed on the bridge. Seed $3000$ at $n=40$ even gets a high cross
accuracy ($0.93$) for the wrong reason (it calls everything dense ``connected''), yet its
discrimination stays at chance ($0.50$) because it then also calls the \emph{split} cliques
connected. This is why we read the conclusion off the mixed model, not this one.}
\label{tab:base_er}
\end{table}

\end{document}
